% Chapter 2 — Architectural Design
\section{Architectural Design}

\subsection{Overview}
% TODO: High-level components and their interaction. One high-level diagram +
%       end-to-end flow: Ingestion/Collector -> message bus -> Storage ->
%       Query/API + Replay -> GUI, plus the Subscription/Notification concern.
%       State tiers and interaction style (event-driven live, request/response
%       historical).

\subsection{Component view}
% TODO: UML component diagram (provided/required interfaces) + a table:
%       component -> responsibility -> key interfaces -> NFR addressed.
%       Fix component names here (they are reused by 2.E and §4).

\subsection{Deployment view}
% TODO: UML deployment diagram + scaling rationale (LB, app servers, DB, broker,
%       cache, GUI). One paragraph per NFR: which decision satisfies PR-1/2/4/7/8.

\subsection{Runtime view}
% TODO: Sequence diagrams describing how components interact to accomplish
%       specific tasks tied to the use cases (e.g. live telemetry push,
%       historical query, replay).

\subsection{Component interfaces}
% TODO: Formal description of each component's interfaces (signatures, data
%       formats, ports). Reuse the exact component names from 2.B.

\subsection{Selected architectural styles and patterns}
% TODO: Per pattern: name -> where used -> which requirement it buys -> why over
%       the alternative. Candidates: Layered, Publish/Subscribe, Modular Monolith
%       vs. Microservices, CQRS, Repository/Adapter (decouples OpenF1).

\subsection{Other design decisions}
% TODO: Decisions not covered above: data model/partitioning, caching, error and
%       retry policies, configuration, logging, security choices.
